\section{Related work}
\label{sec:related}

\subsection{Public automated-driving incident reports}

Public automated-driving crash and disengagement reports have become an important empirical basis for studying automated-vehicle safety outside simulators and proving grounds. The NHTSA Standing General Order database and the California DMV collision and disengagement reports provide rare public evidence about crashes, manual interventions, road contexts, and manufacturer-reported narratives during automated-driving operation \citep{nhtsa2021sgo,californiadmv_collision,californiadmv_disengagement}. Early studies using California reports showed that autonomous-vehicle crashes were often low severity, frequently involved rear-end conflicts, and raised distinct questions about automated-mode operation and control transitions \citep{favaro2017examining,wang2019mechanism,boggs2020crashes}. Disengagement-focused work further characterized who initiated disengagements, what factors were reported, where events occurred, and how system maturity related to disengagement patterns \citep{boggs2020disengagement}. Other studies compared automated-vehicle struck-from-behind crash rates with naturalistic-driving baselines, examined sequence-of-events patterns, and curated crash/disengagement datasets for broader research use \citep{goodall2021struck,song2021sequences,sinha2021data}.

This literature demonstrates the value of public incident reports, but it also motivates the evidence boundary addressed in this paper. Most report-based studies extract patterns, scenario classes, crash partners, movement types, or severity predictors. Recent work has moved closer to narrative analysis by applying topic modeling, text analytics, explainable machine learning, and pre-crash scenario analysis to automated-vehicle crash narratives \citep{alambeigi2020themes,li2024latent,li2025precrash}. These methods can reveal recurring themes and statistically associated factors, but they do not explicitly represent the driver as a controller whose process model is formed and updated through feedback. Consequently, they are not designed to distinguish a report-supported driver-centered pathway from an outcome-driven claim about a supposed takeover failure.

\subsection{Driver takeover, HMI, and process-model evidence}

Human-factors research provides the conceptual reason why driver-side evidence matters. Situation awareness theory emphasizes perception, comprehension, and projection in dynamic systems \citep{endsley1995situation}. Automation research has long shown that use, misuse, disuse, trust, and mode awareness can shape human-automation interaction \citep{parasuraman1997humans,lee2004trust,sarter1995mode}. In automated driving, takeover studies show that time budget, traffic complexity, driver engagement, non-driving-related tasks, and alert modality can affect takeover time and takeover quality \citep{gold2013takeover,eriksson2017takeover,zeeb2015takeover}. Meta-analytic and systematic reviews further document the importance of urgency, non-driving task type, visual versus multimodal alerts, and takeover-request design in helping drivers resume control \citep{zhang2019takeovermeta,salubre2021tor}.

However, the same literature also clarifies why public incident reports are difficult to use for driver-centered claims. Controlled studies can measure driver gaze, reaction time, workload, trust, interface exposure, or response quality. Public crash and disengagement reports usually do not contain those measurements. Treating missing HMI or driver-state evidence as if it were observed would therefore import experimental assumptions into a report setting where they are not warranted. Our method uses this literature to define what kinds of evidence would be needed for a stronger driver-process-model claim, while explicitly blocking claims that the report does not support.

\subsection{LLM and VLM methods for traffic safety analysis}

Recent AAP studies show that large language models and vision-language models can support traffic-safety analysis when embedded in structured workflows. Zhang et al. proposed DDLM, which combines a visual LLM, whole-body pose estimation, and reasoning chains for driver distraction classification and risk assessment \citep{zhang2024ddlm}. Zhang et al. later proposed SeeUnsafe, a multimodal LLM framework for video-based traffic accident analysis, using severity-based aggregation, multimodal prompting, visual grounding, and an information-matching score to analyze accident videos \citep{zhang2025seeunsafe}. Wu et al. integrated crash narratives, diagrams, and street-view imagery with vision-language models to infer crash mechanisms and diagnose roadway risk \citep{wu2026saferroads}.

These studies are important because they show that LLMs and VLMs can be useful for road safety when the task, evidence sources, and evaluation targets are well specified. At the same time, their primary outputs differ from ours. DDLM focuses on observable driver behavior classification; SeeUnsafe focuses on video accident classification and visual grounding; and multimodal roadway-diagnosis work focuses on crash mechanism attribution and infrastructure recommendations. Our focus is narrower and more conservative: given sparse automated-driving incident text, we use an LLM inside an STPA-HF-constrained replay pipeline to generate auditable driver-process-model pathways, blocked claims, and minimum evidence requirements. The richer-evidence case study therefore does not turn the paper into a video-classification study. It tests whether the same replay structure becomes more specific when additional temporal, HMI, driver-action, or scene-evolution evidence reduces uncertainty.

\subsection{STPA-HF and evidence-bounded driver-process-model replay}

STPA models safety as a control problem in which accidents can emerge from unsafe interactions among system components, inadequate feedback, flawed process models, or unsafe control actions \citep{leveson2011engineering,leveson2018handbook}. In human-automation settings, mode confusion research further emphasizes that operators may act unsafely when their mental or process models diverge from the actual system state or automation mode \citep{sarter1995mode,bishop2023modeconfusion}. STPA-HF extends this control-theoretic perspective toward human-factors analysis by asking how feedback, process-model variables, and contextual factors shape human control actions.

Our work builds on this line of research but adapts it to a different evidence regime. Traditional STPA and STPA-HF are often used prospectively for design-time hazard analysis, requirement generation, and scenario analysis. Public automated-driving incident reports, by contrast, are retrospective, sparse, and unevenly documented. We therefore do not claim to perform complete accident reconstruction or complete STPA-HF modeling. Instead, we operationalize a report-compatible subset: source-supported evidence is mapped into driver process-model variables; reported and missing update sources are analyzed; candidate driver actions are generated; unsafe-control-action pathways are derived forward from process-model and action candidates; and outcome information is used only as a compatibility check. This adaptation is the methodological gap addressed by Evidence-Bounded Driver Process-Model Replay.
